% ============================================================================
\section{Discusión, limitaciones y trabajo futuro}
\label{sec:discusion}
% ============================================================================

\subsection{Síntesis}

El núcleo de Tukit encadena tres familias de algoritmos de naturaleza muy
distinta: la lectura con un modelo multimodal, envuelta en reglas
deterministas de aceptación; la clasificación por reglas léxicas con un canal
semántico subordinado; y la optimización combinatoria multiobjetivo. Lo que
las une es una misma postura frente al error. Cada etapa tiene una salida
explícita para «no sé» o «no tengo», y en cada etapa el error barato
(rechazar, abstenerse, mandar a revisión) se prefiere al error caro (aceptar
una lista truncada, ofrecer un producto equivocado, recomendar un kit que no
cabe).

El resultado analítico más importante de este documento es estructural
(\cref{sec:opt-estructura}). Con la formulación vigente, los objetivos son
separables por línea, los extremos se resuelven de forma voraz y exacta, el
compromiso con presupuesto es una mochila de elección múltiple con valores
pequeños y, por lo tanto, polinomial, y el frente completo de cuatro objetivos
se calcula exactamente en $O(n^4k)$. La evaluación empírica es coherente con
ese análisis: el NSGA-II es exacto en los extremos, donde no hace falta, y
aproximado en el frente, donde hay una alternativa exacta.

\subsection{Limitaciones}

\begin{itemize}
  \item \textbf{Sin acoplamiento en el objetivo.} El costo no ve la
        consolidación entre hermanos ni la agrupación de herramientas, que se
        aplican después de elegir, ni un costo por tienda. El optimizador
        resuelve un problema más simple que el real en hogares de varios hijos.
  \item \textbf{Frente aproximado.} El NSGA-II con operadores ordinales sobre
        una variable nominal no garantiza el frente exacto, aunque la siembra
        garantiza no quedar por debajo de la heurística ni de la escalera
        (\cref{prop:siembra}).
  \item \textbf{Umbral de rescate y revisión:} un rescate con coseno de al
        menos 0,93 se acepta sin revisión humana, contra la intención declarada
        del piso de rescate. Su impacto medido es bajo (\cref{tab:cosenos}).
  \item \textbf{Recuperación limitada a la estructura léxica.} Sin el grafo de
        sustituibilidad y sin reordenador, el \emph{recall} de la resolución
        está acotado por los sinónimos escritos y por el rescate semántico.
  \item \textbf{Pesos por omisión} distintos de los medidos en la encuesta
        (\cref{sec:opt-recomendar}).
  \item \textbf{Percepción secuencial}: el tiempo de la primera lectura crece
        linealmente con el número de páginas.
\end{itemize}

\subsection{Trabajo futuro, en orden de impacto}

\begin{enumerate}
  \item \textbf{Frente exacto por programación dinámica.} Reemplazar el
        NSGA-II, o complementarlo como verificación, con la recurrencia del
        \cref{teo:frente-poli}. Para una lista típica son unos $10^4$ estados
        frente a unas $10^7$ operaciones del NSGA-II, y el frente pasaría a ser
        exacto y determinista sin depender de una semilla.
  \item \textbf{Llevar el acoplamiento al objetivo}, si se quiere que el
        optimizador resuelva el problema del hogar completo: la consolidación
        dentro de $f_1$ ---el costo de un producto depende de la demanda
        agregada de los hermanos que lo comparten--- y, eventualmente, un costo
        fijo por tienda. Con ese acoplamiento, la \cref{prop:separable} deja de
        valer y una metaheurística vuelve a estar justificada. En ese caso
        convendría cambiar SBX y PM por operadores categóricos (cruce uniforme
        y mutación por reemplazo aleatorio del índice) y medir con
        hipervolumen \cite{zitzler1999} contra el frente exacto en instancias
        pequeñas. El diseño original preveía también una prueba de falsación:
        sin acoplamiento, la heurística y el NSGA-II deben empatar en
        hipervolumen. Si empatan, queda demostrado que el acoplamiento es la
        contribución.
  \item \textbf{Rescate siempre revisado}: que la respuesta del emparejamiento
        indique que el concepto vino de un rescate semántico y que ese renglón
        vaya a revisión sin importar su puntaje. Requiere agregar un campo al
        contrato entre servicios.
  \item \textbf{Grafo de sustituibilidad y reordenador}, con una evaluación
        que separe «no lo tenemos» de «no lo encontramos» y mida el
        \emph{recall} antes que el ordenamiento.
  \item \textbf{Calibrar} $\tau$ y los pesos de recomendación con datos de uso:
        el umbral de autoaceptación contra las correcciones de los revisores y
        los pesos contra la encuesta.
  \item \textbf{Paralelizar la percepción} por página respetando la cuota del
        proveedor.
  \item \textbf{Evaluar hogares de varios hijos}, que es donde la
        consolidación puede cambiar qué kit es mejor.
\end{enumerate}

\subsection{Conclusión}

Los algoritmos del núcleo son, en su mayoría, sencillos: listas de
decisión, distancias de edición acotadas, coeficientes de Jaccard, centroides,
ordenamientos y programación dinámica. Su valor está en la precisión con que
se delimita cada problema ---qué se decide, con qué información y qué se hace
cuando no se puede decidir---, en que cada regla se justifica con un caso real
y en que las propiedades importantes (no truncar, no fabricar, no restar, no
encarecer, no quedar por debajo de la semilla) se pueden enunciar y demostrar.
El componente más sofisticado, el NSGA-II, es también el que este análisis
señala como el más fácil de reemplazar por algo exacto, mientras la
formulación siga sin acoplamiento entre líneas.
